\documentclass[11pt]{article}
\usepackage{graphicx}
\usepackage{amssymb}
\usepackage{bm}
\usepackage{mathtools}
\usepackage{amsmath}
\usepackage{commath}
\usepackage{amsthm}
\usepackage{enumitem}
\usepackage{float} 
\usepackage{array}
\usepackage{outlines}
\usepackage{setspace}
\usepackage[colorinlistoftodos]{todonotes}
\setlist[enumerate,1]{label=\alph*)}
\setlist[enumerate,2]{label=\roman*)}
\newtheorem{theorem}{Theorem}
\newtheorem{lemma}{Lemma}
\newtheorem{remark}{Remark}
\newtheorem{definition}{Definition}
\newtheorem{prop}{Proposition}
\newcommand*\diff{\mathop{}\!\mathrm{d}}
\usepackage{tikz}
\usetikzlibrary{arrows,shapes,trees, calc}

%%% some of the changes
%%% introduction of c = (c_h, c_3) Coriolis term in the equation for better readability in the estimate parts
%%% the Earth rotation vector is noted now by omega, instead of w // the u,v letters are associated with fluid velocity here, and omega is probably a more classical choice anyway
%%% to have a consistent horizontal component notation, let's choose u_h, with a small h. Let's use H letter only for Sobolev spaces
%%% in the estimates part let's use \partial_3 instead of \partial_z etc since at that point we are after the rescaling process

%%% *** PREAMBLE *** %%%
\title{A weak convergence result for the mathematical equations of the polluted atmosphere expressed in downwind matching coordinates}

\date{}

%%%***  CONTENT *** %%%
\begin{document}

\maketitle

\begin{abstract}

A widely used approach to mathematically describe the atmosphere is to consider it as a geophysical fluid in a shallow domain --- and thus to model it using classical fluid dynamical equations combined with the explicit inclusion of an $\epsilon$ parameter representing the small aspect ratio of the physical domain. In our previous paper we proved a weak convergence theorem for the polluted atmosphere described by the Navier-Stokes equations extended by an advection-diffusion equation. We obtained a justification of the generalised hydrostatic limit model including the pollution effect described in the case of a classical, east-north-upwards oriented local Cartesian coordinates. Here we give an improvement of this statement: we consider a meteorologically more meaningful coordinate system, incorporate the analytical consequences of this coordinate change into the governing equations, and verify that the weak convergence still holds for this altered concentration -- velocity system.

\end{abstract}

{\bf Keywords:} downwind-matching coordinate system, Navier-Stokes equations, shallow domains, pollution evolution equation, hydrostatic approximation, compactness, weak solutions.

\newpage

\tableofcontents
\newpage

%%% *** INTRODUCTION *** %%%
\section{Introduction} \label{Introduction}

A local domain on the surface of the rotating Earth is often described by a Cartesian coordinate system where the $x, y, z$ axes represent the directions towards east, north, and upwards, respectively, as in this specific system the rotation vector, and as a consequence, the Coriolis force takes a simpler form. However, for modelling wind and pollution effects in the atmosphere there is another naturally distinctive viewpoint for fixing the orientation of the coordinate system: the introduction of downwind and crosswind, and to match the horizontal axes according to these meteorological parameters (see \cite{goyal2011mathematical}). In this paper we take this idea as a foundation, and we migrate the result of our previously stated weak convergence theorem to this more considerately chosen coordinate system.

%% Here \todo{write this part} talk a little bit about the original convergence theorem to make the new one more self-contained and for details refer to the first article.

As in this scenario the advection is dominant compared to diffusion effects along the $x$ axis, this approach naturally brings an additional $\epsilon$ term in the Laplacian of the concentration equation, the simplification of the limit model, and brings the challenge of proving a weak convergence result for the product of the vertical velocity and the concentration, as in the energy inequality the $\norm{\partial_1 C^\epsilon}$ term is replaced by $\norm{\sqrt{\epsilon} \partial_1 C^\epsilon}.$ This issue is overcome by using a more regular initial horizontal velocity function, namely $\mathbf{v} \in H^1.$ This additional, stronger initial regularity will result in a smoother $\mathbf{u}^\epsilon$ function, and thus will make up for the lost $H^1$ regularity of the concentration function $C.$

We will also use slightly different, stricter boundary conditions --- these additional requirements are motivated by \textit{a)} the presence of $\epsilon$ next to the sediment functions, and \textit{b)} the need to eliminate certain boundary terms in the new part of the proof, which is based on a result obtained for a periodic domain (\cite{li2019primitive}).

We end up this section with some notations we will use through the paper.

\subsection{Notations}

\begin{enumerate}

\item $\epsilon = $ the aspect ratio,
\item $\nu = $ viscosity,
\item $\Omega_\epsilon, \Omega = $ the original and the rescaled ($\epsilon$-independent) domain, respectively,
\item $\nabla$ stands for the gradient vector; specifically, $(\partial_x, \partial_y, \partial_z)$ before, and $(\partial_1, \partial_2, \partial_3)$ after rescaling,
\item $\nabla_\nu = (\nu_x^{1/2} \partial_x, \nu_y^{1/2} \partial_y, \nu_z^{1/2} \partial_z)$ on $\Omega_\epsilon$, $\nabla_\nu = (\nu_1^{1/2} \partial_1, \nu_2^{1/2} \partial_2, \nu_3^{1/2} \partial_3)$ on the rescaled domain $\Omega,$
\item $ \Delta_\nu$ stands for the anisotropic Laplacian operators $\nu_x \partial^2_{xx} + \nu_y \partial^2_{yy}  + \nu_z \partial^2_{zz}$ and $\nu_1 \partial^2_{11} + \nu_2 \partial^2_{22}  + \nu_3 \partial^2_{33}$ on the original and rescaled domains, respectively,
\item $\mathbf{g}$ represents the force due to gravity, which also includes the centrifugal effect,
\item $\varphi = $ the gravity potential term, i.e. $\mathbf{g} = \nabla \varphi,$
\item $\mathbf{\omega}$ represents the Earth rotation angular speed,
\item $\phi = l(y) = $ latitude, $f = $ the module of the Earth rotation vector
\item $\alpha = -2 f \sin(\theta)\cos(\phi), \beta = 2 f \cos(\theta) \cos(\phi), \gamma = 2 f \sin(\phi),$
\item $\mathbf{c}(\mathbf{u}) = (\mathbf{c}_h (\mathbf{u}), c_3 (\mathbf{u}_h)) = ( -\gamma u_2^\epsilon + \epsilon \beta u_3^\epsilon, \alpha \epsilon u^\epsilon_3 - \gamma u_1^\epsilon, \alpha u^\epsilon_2 - \beta u_1^\epsilon),$ the Coriolis force terms 
\item $\mathbf{K}$ represents the diffusion matrix,
\item $( \cdot,\cdot) = $ the scalar product in $L^2(\Omega)^d$ or the duality $L^p (\Omega)$, $L^{p'} (\Omega),$
\item $< \cdot,\cdot>_{\Gamma} = $ the scalar product or duality on the boundary curve $\Gamma$ ($H^{-1/2}(\Gamma)$, $H^{1/2}(\Gamma)$),
\item $\mathbf{u}_h$ = the horizontal velocity components $(u_1, u_2),$
\item $b(\mathbf{u}_h)$ = $- \gamma (u_2 , u_1),$
\item $L^p_t L^q_x$ is the abbreviated notation for $L^p (0, T; L^q(\Omega)).$

\end{enumerate}

%%% *** THE POLLUTED ATMOSPHERE MODEL *** %%%
\section{The polluted atmosphere model} \label{eqsdescribingatmpluspollution}

As we highlighted before, the key point of this new model is to integrate some of the main ideas of a Gaussian plume dispersion model into the system describing the polluted atmosphere. Therefore we begin by fixing the local Cartesian coordinate system to be adjusted to the downwind direction, i.e. $z$ is oriented upwards, while $x$ is the downwind, and $y$ is the crosswind direction. Another assumption the Gaussian dispersion models typically use as well is that advection is more dominant than diffusion along the downwind direction, which leads to the appearance of an additional $\epsilon$ term in the concentration equation's Laplacian. Namely, we have

\begin{equation} \abs{u_1 \partial_{1} C} \gg \abs{ K_x \partial_{11} ^2 C }\end{equation}

which leads us to using this additional scaling equation

\begin{equation} \label{ksepsk1} K_x = \epsilon K_1.\end{equation}

\begin{remark} \normalfont
  The inclusion of these two features in a pollution model are legitimate when the following circumstances hold:
  \begin{itemize}
    \item We are considering a time interval which is not too long in the sense that it is reasonable to assume that we have a permanent, relatively stable main wind direction, i.e. fixing the horizontal coordinates is reasonable in the first place,
    \item There is an adequate, relatively high wind speed --- low wind speeds in general lead to a situation where three dimensional diffusion dominates, leading to an error caused by the dominant advection assumption. Fast changing wind direction and low windspeed resulting in fast pollutant settling can make it extremely challenging to obtain reliable results from Gaussian models.
  \end{itemize}
\end{remark}

Choosing the coordinate system in this particular way results in changes in three main points of the system. One of these is the Coriolis term, since now the $x$ axis is not necessarily perpendicular to the rotation vector. We need to revisit the boundary conditions as well, and as discussed above, a modification in the pollution concentration equation also has to be taken into account. We will describe these changes and their consequences in detail --- the main challenge in proving this paper's convergence result is in fact to handle the effect of this latter change on the energy inequality.

In order to obtain the set of equations that describe the phenomena we start with the original set of equations

\begin{equation} \label{originalEQ-V}
  \partial_t \mathbf{v} + (\mathbf{v} \cdot \nabla)\mathbf{v} -\Delta_{\nu} \mathbf{v} + \nabla q + 2 \mathbf{\omega} \times \mathbf{v}= \mathbf{g} \text{ in } \Omega_{\epsilon} \times (0,T)
\end{equation}
    
\begin{equation} \label{originalEQ-I}
  \nabla \cdot \mathbf{v} = 0 \text{ in } \Omega_{\epsilon} \times (0,T)
\end{equation}
    
\begin{equation} \label{diffusive}
  \partial_t P + \mathbf{v} \cdot \nabla P = \nabla \cdot (\mathbf{K} \nabla P) + Q,
\end{equation}

As for the rescaling process, here we perform the same steps, with the addition of $(\ref{ksepsk1}).$

To expand the $2 \mathbf{\omega} \times \mathbf{v}$ Coriolis term firstly it is necessary to understand the structure of the new rotation vector. Note that the downwind-matching coordinate system and the original east-north-upwards oriented system are different only in a rotation around the $z$ axis. Let $\theta$ denote the angle between the respective $x$ axes, $\phi$ denotes the latitude, then the rotation vector expressed in the downwind matching coordinate system becomes

\begin{gather}
 \begin{bmatrix} \cos(\theta) & -\sin(\theta) & 0 \\ \sin(\theta) & \cos(\theta) & 0 \\ 0 & 0 & 1 \end{bmatrix} \begin{bmatrix} 0 \\ f \cos(\phi) \\ f \sin(\phi) \end{bmatrix}
 =
  f \begin{bmatrix} -\sin(\theta)\cos(\phi) \\ \cos(\theta) \cos(\phi) \\ \sin(\phi) \end{bmatrix}.
\end{gather}

Now, let

$$\alpha = -2 f \sin(\theta)\cos(\phi), \quad \beta = 2 f \cos(\theta) \cos(\phi), \quad \gamma = 2 f \sin(\phi),$$

and using this, the Coriolis force takes the form

\begin{gather}
\begin{bmatrix} \alpha \\ \beta \\ \gamma \end{bmatrix}
 \times
 \begin{bmatrix} u_1 \\ u_2 \\ \epsilon u_3 \end{bmatrix}
 =
 \begin{bmatrix}  \beta \epsilon u_3 - \gamma u_2 \\ \alpha \epsilon u_3 - \gamma u_1 \\ \alpha u_2 - \beta u_1 \end{bmatrix}.
\end{gather}

Thus, these are the terms that appear in the three velocity equations; and finally, note that the third component is present in the rescaled version of the vertical velocity equation multiplied by $\epsilon,$ which is obtained from the $\epsilon$ term in the denominator of $\partial_z p.$

\begin{remark} \normalfont
The intuitive meaning of what we understand by downwind and the velocity in this new coordinate system: the downwind direction is an approximately steady-in-time direction, it can be seen as the direction marked by the wind direction trackers on an airport or meteorological point. When we consider the downwind, we neglect the small perturbations in its orientation, it can be viewed as a time-averaged direction as long as the fluctuation is relatively small. On the other hand, the velocity is a three dimensional, time dependent vector that changes second after second. 

\end{remark}
    
    We arrive to the merged anisotropic equations of the polluted atmosphere above inland surface:
  
    \begin{equation} \label{ANSC1}
    \partial_t u^\epsilon _1 + \mathbf{u}^\epsilon \cdot \nabla u^\epsilon _1 -\Delta_\nu u^\epsilon _1 -\gamma u_2^\epsilon + \epsilon \beta u_3^\epsilon + \partial_1 p^\epsilon = 0 \text{ in } \Omega \times (0,T)
    \end{equation}
    
    \begin{equation} \label{ANSC2}
    \partial_t u_2^\epsilon + \mathbf{u}^\epsilon \cdot \nabla u_2^\epsilon -\Delta_\nu u_2^\epsilon + \alpha \epsilon u^\epsilon_3 - \gamma u_1^\epsilon + \partial_2 p^\epsilon = 0 \text{ in } \Omega \times (0,T)
    \end{equation}
    
    \begin{equation} \label{ANSC3}
    \epsilon^2 \{  \partial_t u_3^\epsilon + \mathbf{u}^\epsilon \cdot \nabla u_3^\epsilon -\Delta_\nu u_3^{\epsilon} \} + \epsilon \alpha u^\epsilon_2 - \epsilon \beta u_1^\epsilon + \partial_3 p^\epsilon = 0 \text{ in } \Omega \times (0,T)
    \end{equation}
    
    \begin{equation} \label{ANSC4}
    \nabla \cdot \mathbf{u}^\epsilon= 0 \text{ in } \Omega \times (0,T)
    \end{equation}
    
    \begin{equation} \label{ANSC5}
    \partial_t C^\epsilon + \mathbf{u^\epsilon} \cdot \nabla C^\epsilon = \epsilon K_1 \partial_{11} ^ 2 C^\epsilon + K_2 \partial_{22} ^ 2 C^\epsilon + K_3 \partial_{33} ^ 2 C^\epsilon + S_\epsilon \text{ in } \Omega \times (0,T)
    \end{equation}

The initial condition for the velocity and the pollution concentration are
    
    \begin{equation} \label{ANSC7}
    \mathbf{u}^\epsilon (\cdot, t = 0) = \mathbf{u}_0, \quad C^\epsilon (\cdot, t=0) = C_0  \quad \text{ in } \Omega.
    \end{equation}
  
  Finally, the boundary conditions to describe the anisotropic system can be summarised as follows:
  
    \begin{equation} \label{ANSC8}
      \begin{split}
        & \mathbf{u}^\epsilon = 0 \text{ on } \Gamma \times (0,T), \\
        &C^\epsilon = 0 \text{ on } ( \Gamma_U \cup \Gamma_L) \times (0,T), \\
        & \partial_2 C^\epsilon = \partial_3 C^\epsilon = 0 \text{ on } \Gamma_G \times (0,T), \quad \partial_1 C^\epsilon \in L^2 L^2 (0,T; \Gamma_G). \\
      \end{split}
    \end{equation}
    
We highlight the difference between the boundary conditions of the original, east-north oriented, and the new, downwind-matching model.

The difference is the kind of boundary conditions that we can "afford" to have on the lower boundary for the concentration comes from the energy inequality. In the east-north-upwards oriented coordinate system we do not have an $\epsilon$ in the diffusion matrix, and thus it is not an issue to have any of the $\partial_i C$ terms to be different from 0, not even for an ocean case where $\partial_3 C$ becomes nonzero as well (they are required however to be bounded of course). On the other hand, the downwind-matching coordinate system brings in an $\epsilon$ term in front of the $K_1$ diffusion, which makes the model more sensitive to assigning boundary conditions for $C.$ This means that we can not have the relatively relaxed, boundedness requirement for any $\partial_i C$ except for $i=1.$ All other $\partial_2 C, \partial_3 C$ derivatives must be zero on the lower boundary, otherwise they would result in an $(\epsilon - \text{const.})$ negative term.

If we assume $\mathbf{u}^\epsilon = O(1)$, then neglecting the $\epsilon^2$ and $\epsilon$ in the anisotropic equations, we formally arrive to the hydrostatic Navier-Stokes equations (i.e., primitive equations) combined with pollution.

    \begin{equation} \label{HNSC1}    
    \partial_t u_1 + \mathbf{u} \cdot \nabla u_1 - \Delta_\nu u_1 -\gamma u_2 + \partial_1 p = 0 \text{ in } \Omega \times (0,T)
    \end{equation}
    
    \begin{equation} \label{HNSC2}   
    \partial_t u_2 + \mathbf{u} \cdot \nabla u_2 - \Delta_\nu u_2 -\gamma u_1 + \partial_2 p = 0 \text{ in } \Omega \times (0,T)
    \end{equation}
    
    \begin{equation} \label{HNSC3}   
    \partial_3 p = 0 \text{ in } \Omega \times (0,T)
    \end{equation}
    
    \begin{equation} \label{HNSC4}   
    \nabla \cdot \mathbf{u} = 0 \text{ in } \Omega \times (0,T)
    \end{equation}
    
    \begin{equation} \label{HNSC5}   
    \partial_t C + \mathbf{u} \cdot \nabla C = K_2 \partial_{22} ^ 2 C + K_3 \partial_{33} ^2 C + S \text{ in } \Omega \times (0,T)
    \end{equation}
    
    \begin{equation} \label{HNSC6}   
    u_1 = u_2 = u_3 n_3 = 0 \text{ on } \Gamma \times (0,T)
    \end{equation}
    
    \begin{equation} \label{HNSC7}   
    u_i (\cdot, t=0) = u_{0 i}, C (\cdot, t=0) = C_0 \text{ in } \Omega, \text{$i$ = 1,2}
    \end{equation}
    
    \begin{equation} \label{HNSC8}
      \begin{split}
        &C = 0 \text{ on } ( \Gamma_U \cup \Gamma_L) \times (0,T), \\
        & \partial_2 C = \partial_3 C = 0 \text{ on } \Gamma_G \times (0,T), \partial_1 C \in L^2 L^2 (0,T; \Gamma_G) \\
      \end{split}
    \end{equation}
    
\subsection{Energy inequality}

Find $\mathbf{u}^\epsilon = (\mathbf{u}^\epsilon_H, u^\epsilon_3) \in L^2 (0,T ; \mathbf{V}) \cap L^\infty (0, T; L^2 (\Omega)^3)$ and $C^\epsilon \in L^\infty (0,T; L^2(\Omega)) \cap L^2 (0,T; H^1(\Omega))$, such that
  
  \begin{equation}
    \begin{split} \nonumber
      & \int_0^T \bigg [ -(\mathbf{u}^\epsilon_H, \partial_t \mathbf{\tilde{u}}_H) + (\nabla_\nu \mathbf{u}^\epsilon_H, \nabla_\nu \mathbf{\tilde{u}}_H) - (\mathbf{u}^\epsilon_H, (\mathbf{u}^\epsilon \cdot \nabla) \mathbf{\tilde{u}}_H) + (b(\mathbf{u}^\epsilon_H), \mathbf{\tilde{u}}_H) \bigg ] \diff t \\
      & \hspace{-0.5cm} + \int_0 ^ T \epsilon \beta \big [ (u_3 ^\epsilon, \tilde{u}_1) - (u_1^\epsilon, \tilde{u}_3) \big ] + \epsilon^2\big [ -(u_3^\epsilon, \partial_t \tilde{u}_3) +(\mathbf{u}^\epsilon \cdot \nabla u_3 ^ \epsilon, \tilde{u}_3) + (\nabla_\nu u_3^\epsilon, \nabla_\nu \tilde{u}_3) \big ] \diff t \\
      & + \int_0 ^ T \epsilon \alpha \big [ (u_3 ^\epsilon, \tilde{u}_2) - (u_2^\epsilon, \tilde{u}_3) \big ] = (\mathbf{u}^\epsilon_{0H}, \mathbf{\tilde{u}}_H(0)) + \epsilon^2 (u_{03}^\epsilon, \tilde{u}_3 (0)) \diff t
    \end{split}
  \end{equation}
  
  and
  
    \begin{equation}
    \begin{split} \nonumber
      & \int_0^T \bigg [ -(C^\epsilon, \partial_t \tilde{C}) -(\mathbf{u}^\epsilon C^\epsilon, \nabla \tilde{C}) + (\mathbf{M} \nabla C^\epsilon, \nabla \tilde{C}) - \int_{\Gamma_G} \boldsymbol{\sigma} \mathbf{M} \tilde{C} \vec{\boldsymbol{n}}_{\Gamma_G} \bigg ] \diff t \\
      & =  (C^\epsilon (0), \tilde{C} (0)) + \int_0^T (S^\epsilon,\tilde{C}) \diff t
    \end{split}
  \end{equation}
  
    \hspace{0.2cm} for all $\mathbf{\tilde{u}} = (\mathbf{\tilde{u}}_H, \tilde{u}_3) \in H^1 (0,T;\mathbf{V})$ with $\mathbf{\tilde{u}}(T) = 0$ and $\tilde{C}$ such that
    
    \hspace{0.2cm} $\tilde{C} \in H^1 (0,T, H^2_{\Gamma_A} (\Omega))$ and $\tilde{C} (T) = 0$.


%%% *** WEAK FORMULATION AND MAIN RESULT*** %%%
\section{Weak formulation and Main Result} \label{weakformulationmainresult}

The main result is analogous to the original one, stated for the case of \todo{where is the weak continuity used in the proof?} Leray-Hopf weak solutions and $\mathbf{v}_0 \in H^1.$ The domain is a bounded Lipschitzian domain.

\begin{theorem} \label{maintheorem}
  Let $\mathbf{v}_0$ $\in$ $H^1 (\Omega)$, ${u_3}_0 \in L^2(\Omega)$ with $\nabla \cdot \mathbf{u}_0 = 0, \mathbf{u}_0 \cdot \mathbf{n} = 0$ on $\partial \Omega$; let $C_0 \in$ $L^2 (\Omega)$, $C_0 = 0$ on $ \partial \Omega$, and assume that the boundary conditions (\ref{ANSC8}) hold; then as the aspect ratio $\epsilon$ tends to zero, any weak solution $(\mathbf{u}^\epsilon, C^\epsilon)$ of the anisotropic equations (\ref{ANSC1}) - (\ref{ANSC5}) converge to a weak solution $(\mathbf{u}, C)$ of the hydrostatic equations of the polluted atmosphere (\ref{HNSC1}) - (\ref{HNSC5}).
\end{theorem}

In the theorem of \cite{li2019primitive} they also use additional requirements such as zero average \todo{why can't they simply use the Poincare inequality and refer to the fact that they don't have boundary values?} for the initial data $\int_\Omega v_0 = 0$.

Note that in their case they work with a \todo{include the Coriolis term} Coriolis-term-free model. However these terms are linear and typically non-problematic.
  
%%% *** PROOF *** %%%
\section{Proof of the main theorem} \label{proof}

The proof is mainly based on the Strong Convergence Justification paper's convergence result. In \cite{li2019primitive} it is proved that a \todo{they use L-H WS, we worked on WS} Leray-Hopf weak solution $(\mathbf{v}_\epsilon, w_\epsilon)$ of (SNS) converges strongly to the strong solution $(\mathbf{v}, w)$ in $L^2 L^2,$ provided that we have $\mathbf{v}_0 \in H^1,$ and the domain is periodic (see Theorem 1.1 on p5).

The strong convergence of $w_\epsilon$ to $w$ allows us to bring the proof to a closure, obtaining that $u_3^\epsilon C^\epsilon$ tends to $u_3 C.$

We verify here that Theorem 1.1 in \cite{li2019primitive} is also valid for the case of a bounded domain that is not periodic (i.e. we have boundary segments). We follow closely the proof provided in Section4 and update each point where the periodicity of the domain is applied in one way or another.

The switch from a periodic domain to the bounded Lipschitz domain effects the proof in the following ways in detail:

\begin{enumerate}[label=\roman*)]
  \item The applicability of the Poincare inequality
  \item Integration by parts, possible additional boundary terms in the weak form, etc.
\end{enumerate}

We will check these steps by step.

Note that in this case we have the more regular $v_0 \in H^1$ condition, so we will use the $u_3 = 0$ equality on the boundary directly, instead of the original $u_3 n_3 = 0.$

\subsection{The applicability of the Poincare inequality}

Concerning the first point, i.e. the applicability of the Poincare inequality, we have to check two main types of scenarios:

\subsubsection{First order terms}
  Estimating the function by its gradient, ie. $\norm{\mathbf{v}}$ by $\norm{\nabla \mathbf{v}}, $ $\norm{\mathbf{v}^\epsilon}$ by $\norm{\nabla \mathbf{v}^\epsilon}, $ $\norm{w}$ by $\norm{\nabla w}, $ and $\norm{w^\epsilon}$ by $\norm{\nabla w^\epsilon}. $
  
  These points are automatically verified by using following version of the Poincare inequality (which is a different variant from the zero spatial average version that is used in the article): for a function in $W^{1,2}_0$ we have that $\norm{u}_2 \leq \norm{\nabla u}_2.$ Here $W^{1,2}_0$ ($k=1, p=2$) stands for $H^1$ functions whose $0, ... , n$-th derivative is zero on the boundary, $n \leq k-1 = 1-1 = 0.$ Since the $\mathbf{u}^\epsilon, \mathbf{u}$ functions are zero in the boundary, this zero-trace version of the Poincare inequality can be applied directly.
  
  This means that $\norm{v}_6 \leq C \norm{\nabla v}_2$ is valid in p21, $\norm{V^\epsilon} \leq C \norm{\nabla V^\epsilon}$ on p22 is correct as well, etc.
 
\subsubsection{Second order terms}

Estimating the gradient by the Laplacian.

\begin{enumerate}

\item p17, updating the inequality which originally had the following form:

$$ \hspace{-1cm} \int_0 ^ T \int_M \int_{-1}^1 \abs{u^\epsilon} \abs{\nabla w^\epsilon} \diff z \int_{-1}^1 \abs{\nabla_H v} \leq C \int_0^T \norm{u^\epsilon}_2^{1/2} \norm{\nabla u^\epsilon}_2^{1/2} \norm{\nabla w^\epsilon}_2 \norm{\nabla_H v}_2^{1/2} \norm{\Delta v}_2^{1/2}.$$

This inequality is a result of Lemma 2.1 ($h= u^\epsilon, g = \nabla w^\epsilon$ and $f = \nabla_H v$) combined with the Poincare inequality.

As we don't have $\nabla v = 0$ on the boundary, we can not use $\norm{f} \leq \norm{\nabla f},$ since we don't have $\norm{\nabla v} \leq \norm{\Delta v}.$ Thus, instead of the original bound we will have

$$C \int_0^T \norm{u^\epsilon}_2^{1/2} \norm{\nabla u^\epsilon}_2^{1/2} \norm{\nabla w^\epsilon}_2 \norm{\nabla_H v}_2^{1/2} (\textcolor{orange}{\norm{\nabla_H v}_2^{1/2}} + \norm{\Delta v}_2^{1/2}).$$

We know from the article that we have a bound for

$$ C \int_0^T \norm{u^\epsilon}_2^{1/2} \norm{\nabla u^\epsilon}_2^{1/2} \norm{\nabla w^\epsilon}_2 \norm{\nabla_H v}_2^{1/2} \norm{\Delta v}_2^{1/2}, $$

on the other hand, for the newly added term we can use

\begin{equation}
  \begin{split} \nonumber
    & C \int_0^T \norm{u^\epsilon}_2^{1/2} \norm{\nabla u^\epsilon}_2^{1/2} \norm{\nabla w^\epsilon}_2 \norm{\nabla_H v}_2 \\
    & \leq C \sup_{0 \leq t \leq T} \norm{u^\epsilon}_2^{1/2} \norm{\nabla v}_2 \bigg( \int_0^T \norm{\nabla u^\epsilon}_2 \diff t \bigg)^{1/2} \bigg ( \int_0^T \norm{\nabla w^\epsilon}_2^2 \diff t \bigg)^{1/2}.
  \end{split}
\end{equation}

So the integral is still valid.

\item Updating the bounds I2 and I3, and with those the final steps of the proof as well.

Firstly, on p22, we update the bound for $I''_2:$

$$ \int_0^{t_0} \int_M \bigg ( \int_{-1}^1 \abs{\nabla_H V^\epsilon} \diff z' \bigg) \bigg ( \int_{-1}^1 \abs{V^\epsilon} \abs{\partial_z v} \bigg ).$$

Again we will use Lemma 2.1 and the Poincare inequality (for $V^\epsilon$ we can still use it). We choose $f = \nabla_H V^\epsilon, g = V^\epsilon, h = \partial_z v.$

The difference this time is that instead of $\norm{h}_2^{1/2} \norm{\nabla h}_2^{1/2}$ we will have the complete $\norm{h}_2^{1/2} (\norm{h}_2^{1/2} + \norm{\nabla h}_2^{1/2}),$ as we can not apply the Poincare inequality for $\nabla v.$ 

This means that eventually we arrive to the bound

\begin{equation}
  \begin{split}
     & C \int_0^{t_0} \norm{\nabla V^\epsilon}_2^{3/2} \norm{V^\epsilon}_2^{1/2} \norm{\nabla v}_2^{1/2} (\textcolor{orange}{\norm{\nabla v}_2^{1/2}} + \norm{\Delta v}_2^{1/2})  \\
     & \leq \frac{1}{18} \norm{\nabla V^\epsilon}^2_2 + C \int_0^{t_0} \norm{\nabla v}_2^2 (\textcolor{orange}{\norm{\nabla v}_2^{1/2}} + \norm{\Delta v}_2^{1/2})^4 \norm{V^\epsilon}_2^2 \diff t \\
     & \leq \frac{1}{18} \norm{\nabla V^\epsilon}^2_2 + \textcolor{orange}{C \int_0^{t_0} \norm{\nabla v}_2^4 \norm{V^\epsilon}_2^2 \diff t} + C \int_0^{t_0} \norm{\nabla v}_2^2  \norm{\Delta v}_2^2 \norm{V^\epsilon}_2^2 \diff t,
  \end{split}
\end{equation}

using the Young inequality.

With this we obtain the updated version of (4.13) as well:

$$ I_2 \leq \frac{1}{6} \norm{\nabla V^\epsilon}^2_2 + \textcolor{orange}{C \int_0^{t_0} \norm{\nabla v}_2^4 \norm{V^\epsilon}_2^2 \diff t} + C \int_0^{t_0} \norm{\nabla v}_2^2  \norm{\Delta v}_2^2 \norm{V^\epsilon}_2^2 \diff t$$

Similarly, at the bound of I3 on p23 we arrive to this extended version with an additional $\nabla v$ term:

\begin{equation}
  \begin{split} \nonumber
    I3 & \leq C \epsilon^2 \int_0 ^ {t_0} ( \norm{v^\epsilon}_2^2 \norm{\nabla v^\epsilon}_2^2 + \norm{\nabla v}_2^2 (\textcolor{orange}{\norm{\nabla v}_2^{1/2}} + \norm{\Delta v}_2^{1/2})^4 + \epsilon^2 \norm{w^\epsilon}_2^2 \norm{\nabla w^\epsilon}_2^2 ) \diff t \\
    & + \frac{1}{6} \big ( \norm{\nabla V^\epsilon}_2^2 + \epsilon^2 \norm{\nabla W^\epsilon}_2^2  \big)
  \end{split}
\end{equation}

After this we arrive back to the original estimate of I3, because the new $\norm{\nabla v}_2^2 \cdot \norm{\nabla v}_2^2$ term can be estimated by the same Corollary 3.1, by the same $C(\norm{v_0}_{H^1}, L_1, L_2)$ term. So even though we had an additional term for some part of the estimate of I3, its estimate melts back into the one we originally had.

Finally, the closure of the proof, following the changes in I2, becomes the following.

\begin{equation}
  \begin{split} \nonumber
  f(t):= & (\norm{V^\epsilon}_2^2 + \epsilon^2 \norm{W^\epsilon}_2^2)(t) + \int_0^t (\norm{\nabla V^\epsilon}_2^2 + \epsilon^2 \norm{\nabla W^\epsilon}_2^2) \diff s \\
  & \leq C \epsilon^2 [(\norm{v_0}_2^2 + \epsilon^2 \norm{w_0}_2^2)^2 + C(\norm{v_0}_{H^1}, L_1, L_2)] \\
  & + \textcolor{orange}{C \int_0^{t} \norm{\nabla v}_2^4 \norm{V^\epsilon}_2^2 \diff s} + C \int_0^{t} \norm{\nabla v}_2^2  \norm{\Delta v}_2^2 \norm{V^\epsilon}_2^2 \diff s =: F(t),
  \end{split}
\end{equation}

With this we have

\begin{equation}
  \begin{split} \nonumber
   F'(t) = & \textcolor{orange}{C \norm{\nabla v}_2^4 \norm{V^\epsilon}_2^2} + C \norm{\nabla v}_2^2  \norm{\Delta v}_2^2 \norm{V^\epsilon}_2^2 \\
   & \leq C (\textcolor{orange}{\norm{\nabla v}_2^4} + \norm{\nabla v}_2^2  \norm{\Delta v}_2^2)f(t) \leq C (\textcolor{orange}{\norm{\nabla v}_2^4} + \norm{\nabla v}_2^2  \norm{\Delta v}_2^2)F(t),
  \end{split}
\end{equation}

from which, applying the Gronwall inequality, and using Corollary 3.1 we deduce

\begin{equation}
  \begin{split} \nonumber
    f(t) \leq & F(t) \leq e^{C \int_0 ^ t (\textcolor{orange}{\norm{\nabla v}_2^4} + \norm{\nabla v}_2^2  \norm{\Delta v}_2^2) \diff s} F(0) \leq e^{C(\norm{v_0}_{H^1}, L_1, L_2)} F(0)\\
    & \leq \epsilon^2 C(\norm{v_0}_{H^1}, L_1, L_2) (\norm{v_0}_2^2 + \epsilon^2 \norm{w_0}_2^2 + 1)^2,
  \end{split}
\end{equation}

which implies the conclusion.

\end{enumerate}

\subsection{Integration by parts}

\begin{itemize}

  \item The weak formulation on p17 --- there are no additional boundary terms in the non-periodic case. In terms of boundary terms there are no additional requirements.
  
  \item Integration by parts on p18
  
    $$<\partial_t w, w^\epsilon> = - \int_{\partial \Omega} \bigg ( \int_0^z \partial_t v \diff z' \bigg ) w^\epsilon n_H + \int_\Omega \bigg ( \int_0^z \partial_t v \diff z' \bigg) \nabla_H \cdot w^\epsilon,$$
    
    from where we arrive easily to 
    
     $$<\partial_t w, w^\epsilon> = \int_\Omega \bigg ( \int_0^z \partial_t v \diff z' \bigg) \nabla_H \cdot w^\epsilon,$$
     
     that was required, by using that $w^\epsilon = 0$ on $\partial \Omega.$
     
   \item Integration by parts on p19
   
   $$<\partial_t w, w^\epsilon - w> = \int_\Omega \bigg ( \int_0^z \partial_t v \diff z' \bigg) \nabla_H \cdot (w^\epsilon - w),$$ as the boundary term $\int_{\partial \Omega} \bigg ( \int_0^z \partial_t v \diff z' \bigg ) (w^\epsilon - w) n_H$ becomes zero, since in this version  - unlike the original one, where we had $u_3 n_3 = 0$ on $\Gamma$ - we also have $w = 0$ on $\Gamma$ (and of course $w^\epsilon = 0$ on the boundary as well). 
   
   \item Verifying (4.8).
   
   Multiplying $\partial_t v + (u \cdot \nabla) v - \Delta v + \nabla_H p = 0$ by $v^\epsilon$ and integrating on $Q_{t_0},$ the boundary term from the Laplacian falls out, as $\int_{\partial \Omega} \nabla v \cdot v^\epsilon = 0$ because  $v^\epsilon = 0$ on the boundary. On the other hand, the pressure terms vanishes since $ \int_\Omega \nabla_H p \cdot v^\epsilon = \int_\Omega (\nabla_H p \cdot v^\epsilon + \partial_z p w^\epsilon),$ using that $\partial_z p = 0,$ and thus
   
   $$ \int_\Omega \nabla_H p v^\epsilon = \int_{\partial \Omega} p \cdot \mathbf{u}^\epsilon \vec{n} - \int_\Omega p \cdot ( \nabla \cdot \mathbf{u}^\epsilon) = 0,$$
   
   using the divergence-free condition and $\mathbf{u}^\epsilon = 0$ on $\partial \Omega.$
   
   \item Verifying (4.9): in terms of boundary conditions the analogous ideas can be applied as in the previous step.
   
   \item On p22, the calculation of the I2 term trivially remains correct as $\int_{\partial \Omega} \mathbf{u} \mathbf{v} \mathbf{v}^\epsilon \vec{n} = 0.$ The calculation of $I''_2$ stays valid as well since
   $\int_{\partial \Omega} W^\epsilon V^\epsilon v n_3 = 0.$

\end{itemize}

\nocite{*}
\bibliographystyle{abbrv}
\bibliography{the_u_c_eps_model_using_strong_conv_u3_SCJ}

\end{document}

%%% add http://elte.prompt.hu/sites/default/files/tananyagok/AtmosphericChemistry/ch10s03.html
%%% http://twister.caps.ou.edu/PM2000/Chapter7.pdf
